\subsubsection{Poisson Equation}

\begin{framed}
\textbf{Interactive Visualization !}\\
This notebook contain interactive visualization that needs binder connection. Please click the \textbf{power button on the top right corner} before you start reading ! Establishing the connection and building the container takes a while, so clicking it before you read would make sure the interactive code to be ready to run when you scroll on to them.
\end{framed}

\paragraph{2D Poisson Equation}

Let $\Omega \subset \mathbb{R}^2$, we consider the Poisson equation with homogenous boundary condition:

\begin{definition}[2D Poisson Equation with Homogenous Boundary Condition]\begin{equation}
\label{2dpoihb}
  -\Delta u &= f(x,y) \qquad (x,y)\in \Omega \\
  u(x,y) &= 0 \qquad (x,y)\in \Gamma_1 \\
  \frac{\partial u}{\partial \nu}(x,y) &= 0 \qquad (x,y)\in \Gamma_1
\end{equation}

\end{definition}where $\Gamma_1\cup \Gamma_2 = \partial \Omega$ such that both $\Gamma_1, \Gamma_2$ is measure-zero. $\nu$ is the outward normal vector.

\paragraph{Spectral Method}

Let
\begin{equation}
u(x,y)=\sum_{-\infty<p,q<\infty} a_{pq}  \exp(i(px +qy))
\end{equation}
Which has:
\begin{equation}
\Delta u = \sum_{-\infty<p,q<\infty} a_{pq}
\end{equation}
Also let the fourier transform of $f(x,y)$ be $\mathcal F (f(x,y))=F(p,q)$.

\paragraph{Finite-Element Method(FEM)}

\subparagraph{The Tent Function}

We consider the

\subparagraph{Viualizing the Tent Basis}

\begin{framed}
\textbf{Tip}\\
If you have not do so, click the power button on the top right corner.
After everything builds up, you can click the start button ▶️. You should see two sliders showing up below.

\begin{itemize}
\item The resolution slider determines how ``dense'' these tent basis are in filling the domain. Lower resolution means larger and less tents and vice versa.
\item The position slider let you adjust which particular basis function you are watching. In otherword, where the ``tent'' is located in the domain $\Omega$
\end{itemize}
\end{framed}

%  We define:
% :::{prf:definition} Dirchlet Integral
% Let $u\in C^2(\Omega), v\in C^{\infty}(\Omega)$, define the bilinear form on $H^1(\Omega)$ to be
% $$
% a(u,v) := \int_{\Omega}\nabla u \nabla v
% $$
% :::

We can show that:

\begin{proposition}Suppose $u \in C^2(\bar \Omega)$ is a solution of (\ref{2dpoihb}), then $\forall v \in H^1_{\Gamma_1}(\Omega)$, we have:
\begin{equation}
\int_{\Omega}\nabla u \nabla v=\int_{\Omega} fv
\end{equation}

\end{proposition}\begin{proof}We first show \textit{Green's Formula}(change of variable):
\begin{equation}
\int_{\Omega}v \Delta u = \int_{\partial \Omega} v \frac{\partial u}{\partial \nu} - \int_{\Omega} \nabla u \nabla v
\end{equation}

\end{proof}This motivates us to define a solution in a weaker sense. Since $a(u,v)=(f,v)$ is a necessary condition of $u$, how about making every $u$ that satisfies this equation a ``weak solution''(or generalized solution)?

We first define the space that our weak solution would reside:
\begin{equation}
H^1_{\Gamma_1}(\Omega)=\{f\in H^1(\Omega): f=0 \text{ on } \Gamma_1 \}
\end{equation}